\section{Conclusion}

The contribution of this work is an audited classification engine and, built on the trust that engine earns, a hierarchy that determines when a measured linewidth may legitimately be read as a mechanical damping ratio. The engine's three core scientific safeguards (exceptional-point-versus-semisimple discrimination by eigenvector deficiency, provenance gating that scopes an exceptional-point verdict to the supplied representation, and a constructive counterexample demonstrating that local spectral architecture alone cannot supply a reaction rate), together with the three implementation validity gates that keep them trustworthy in practice (scalar-\(\chi\) applicability, source gating of numerical inputs, and finite-precision equality caution), answer questions a published damped-oscillator fit cannot. Around that engine, the damping ratio \(\chi = \gamma/2\omega_0\) supplies a coordinate that reads the same oscillatory-to-monotone boundary wherever a second-order or decoupled modal representation is established, with a qualifying rule that admits only amplitude damping to that coordinate.

Part~III sharpens that qualifying rule into a decision procedure. A measured spectral line is not self-describing: it may combine static inhomogeneous broadening with a homogeneous width, and the homogeneous width may itself contain lifetime relaxation, pure dephasing and an orientational contribution, of which only the lifetime part is a mechanical damping ratio. A three-step procedure separates them: homogeneous isolation; then separation of the lifetime, pure-dephasing, and orientational contributions, requiring an independent lifetime measurement and an orientational term that is measured, bounded, or independently excluded; then mechanical placement, so that a \(\chi\) is computed only from a lifetime certified as such. The ratio \(R_\phi = (1/T_\phi)/(1/2T_1)\) quantifies the bias (\(1 + R_\phi + R_{\mathrm{or}}\)) incurred by reading the total homogeneous width as lifetime damping, with \(R_\phi = 1\) the equal-contribution crossover rather than a boundary of recoverability. The procedure is carried through one worked real system, the CO stretch of Rh(CO)$_2$acac, anchored to the two endpoints the source reports: the homogeneous width grows by a factor of about 13.6 between them (0.11 to 1.5~cm$^{-1}$), and the source establishes that pure dephasing becomes dominant, since the pump-probe lifetime varies far more weakly than the echo coherence time. A source-informed linear reading of the plotted lifetime data places the mechanical damping ratio near $3\times10^{-5}$ at both endpoints, but that numerical value is illustrative until the published lifetime points are calibrated and digitized; the quantitative claim is limited to the width growth, so coherence loss and mechanical damping are demonstrably different in kind even where the damping value itself remains illustrative. The parallel with the engine is exact, in that a degeneracy cannot be classified from eigenvalues alone and a linewidth cannot be classified from its total width alone, and in both cases the honest object is a test rather than a table. The parallel is also computational and not merely thematic: the fitted lifetime at each temperature builds a scalar generator that the Part~I classifier reads back as a stable oscillatory simple mode whose damping coordinate matches the directly computed $\chi$, so the linewidth procedure and the degeneracy engine are demonstrably one pipeline.

\section*{Data and code availability}

The numerical engine requires NumPy and SciPy; figure generation uses Matplotlib and the regression suites use pytest, with tested versions recorded in the release manifest and printed by the package validator. The code runs on a laptop and ships with a reproducibility guide and adversarial test suites. A master validator (\texttt{run\_all\_chemsa.py}) is included; the validator treats a missing file as a failure and requires an explicit success signal from every stage, and reports twenty of twenty stages passing when \texttt{pytest} is available. A separate package validator (\texttt{validate\_package.py}) checks the release as a whole in ten stages, including that every shipped file except \texttt{MANIFEST.json} itself carries a role and a SHA-256 hash in that manifest (the archive checksum authenticates the manifest), and that no undeclared file (a stray runtime artifact, for instance) is present in the release. The biorthogonal response layer is \texttt{biorthogonal\_response.py} with \texttt{test\_biorthogonal.py}; with the crowded-spectrum certification suite (\texttt{test\_crowded\_spectrum.py}) and the provenance source-gate suite (\texttt{test\_provenance\_gate.py}), the engine and safeguard suites comprise 413 tests, all passing. Randomized statistics reproduce from \texttt{test\_randomized.py} with a fixed seed, the external exceptional-point cross-check from \texttt{benchmark\_ep\_pendulum.py}, and the over-damped CsPbBr\(_3\) cell from \texttt{verify\_cspbbr3\_cell.py}. The source gate (\texttt{provenance\_gate.py}) makes an un-sourced number un-usable and is exercised on the worked example: the reported endpoint width growth is computed through the gate from the two source-anchored widths, so an illustrative point cannot enter it. The single worked example of Part~III and its figure reproduce from \texttt{linewidth\_hierarchy.py}, which reads points digitized from the cited source's Figure~2 (shipped as \texttt{rhco2acac\_fig2\_raw\_digitized.csv}), fits the source-described linear $2T_1(T)$ relation, and derives the decomposition and $\chi$. Every internally generated numerical result and every reported source-derived transformation reproduces from the provided scripts and declared input files; none is asserted from memory. Several application scripts begin from literature inputs transcribed by hand into the shipped data files, and those transcriptions are the declared inputs rather than script-generated quantities.
